% problem_id: S001-lemma-desingularize
% source_id: S001 (Stacks Project)
% source_file: smoothing.tex
% statement_locator: smoothing.tex lines 1664--1679
% proof_locator: smoothing.tex lines 1681--1809
% env: lemma
% id_note: label 在本来源内唯一
% pairing: tight (verified)
% status: complete (statement + author proof verbatim + reason + locators)
%
\begin{lemma}
\label{lemma-desingularize}
Let $R$ be a Noetherian ring.
Let $\Lambda$ be an $R$-algebra. Let $\pi \in R$ and
assume that $\text{Ann}_\Lambda(\pi) = \text{Ann}_\Lambda(\pi^2)$. Let
$A \to \Lambda$ be an $R$-algebra map with $A$ of finite
presentation. Assume
\begin{enumerate}
\item the image of $\pi$ is strictly standard in $A$ over $R$, and
\item there exists a section $\rho : A/\pi^4 A \to R/\pi^4 R$
which is compatible with the map to $\Lambda/\pi^4 \Lambda$.
\end{enumerate}
Then we can find $R$-algebra maps $A \to B \to \Lambda$ with $B$
of finite presentation such that $\mathfrak a B \subset H_{B/R}$ where
$\mathfrak a = \text{Ann}_R(\text{Ann}_R(\pi^2)/\text{Ann}_R(\pi))$.
\end{lemma}

\begin{proof}
Choose a presentation
$$
A = R[x_1, \ldots, x_n]/(f_1, \ldots, f_m)
$$
and $0 \leq c \leq \min(n, m)$ such that
(\ref{equation-strictly-standard-one}) holds for $\pi$ and such that
\begin{equation}
\label{equation-star}
\pi f_{c + j} \in (f_1, \ldots, f_c) + (f_1, \ldots, f_m)^2
\end{equation}
for $j = 1, \ldots, m - c$. Say $\rho$ maps $x_i$ to the class of
$r_i \in R$. Then we can replace $x_i$ by $x_i - r_i$. Hence we may
assume $\rho(x_i) = 0$ in $R/\pi^4 R$. This implies that
$f_j(0) \in \pi^4R$ and that $A \to \Lambda$ maps $x_i$
to $\pi^4\lambda_i$ for some $\lambda_i \in \Lambda$. Write
$$
f_j = f_j(0) + \sum\nolimits_{i = 1, \ldots, n} r_{ji} x_i + \text{h.o.t.}
$$
This implies that the constant term of $\partial f_j/\partial x_i$ is
$r_{ji}$. Apply $\rho$ to (\ref{equation-strictly-standard-one})
for $\pi$ and we see that
$$
\pi = \sum\nolimits_{I \subset \{1, \ldots, n\},\ |I| = c}
r_I \det(r_{ji})_{j = 1, \ldots, c,\ i \in I} \bmod \pi^4R
$$
for some $r_I \in R$. Thus we have
$$
u\pi = \sum\nolimits_{I \subset \{1, \ldots, n\},\ |I| = c}
r_I \det(r_{ji})_{j = 1, \ldots, c,\ i \in I}
$$
for some $u \in 1 + \pi^3R$. By
Algebra, Lemma \ref{algebra-lemma-matrix-left-inverse}
this implies there exists a $n \times c$ matrix $(s_{ik})$ such that
$$
u\pi \delta_{jk} = \sum\nolimits_{i = 1, \ldots, n} r_{ji}s_{ik}\quad
\text{for all } j, k = 1, \ldots, c
$$
(Kronecker delta). We introduce auxiliary variables
$v_1, \ldots, v_c, w_1, \ldots, w_n$ and we set
$$
h_i = x_i - \pi^2 \sum\nolimits_{j = 1, \ldots c} s_{ij} v_j - \pi^3 w_i
$$
In the following we will use that
$$
R[x_1, \ldots, x_n, v_1, \ldots, v_c, w_1, \ldots, w_n]/
(h_1, \ldots, h_n) = R[v_1, \ldots, v_c, w_1, \ldots, w_n]
$$
without further mention. In
$R[x_1, \ldots, x_n, v_1, \ldots, v_c, w_1, \ldots, w_n]/
(h_1, \ldots, h_n)$ we have
\begin{align*}
f_j & = f_j(x_1 - h_1, \ldots, x_n - h_n) \\
& =
\pi^2 \sum\nolimits_{k = 1}^c
\left(\sum\nolimits_{i = 1}^n r_{ji} s_{ik}\right) v_k
+
\pi^3 \sum\nolimits_{i = 1}^n r_{ji}w_i \bmod \pi^4 \\
& =
\pi^3 v_j + \pi^3 \sum\nolimits_{i = 1}^n r_{ji}w_i \bmod \pi^4
\end{align*}
for $1 \leq j \leq c$. Hence we can choose elements
$g_j \in R[v_1, \ldots, v_c, w_1, \ldots, w_n]$
such that $g_j = v_j + \sum r_{ji}w_i \bmod \pi$
and such that $f_j = \pi^3 g_j$ in the $R$-algebra
$R[x_1, \ldots, x_n, v_1, \ldots, v_c, w_1, \ldots, w_n]/
(h_1, \ldots, h_n)$. We set
$$
B = R[x_1, \ldots, x_n, v_1, \ldots, v_c, w_1, \ldots, w_n]/
(f_1, \ldots, f_m, h_1, \ldots, h_n, g_1, \ldots, g_c).
$$
The map $A \to B$ is clear. We define $B \to \Lambda$ by mapping
$x_i \to \pi^4\lambda_i$, $v_i \mapsto 0$, and $w_i \mapsto \pi \lambda_i$.
Then it is clear that the elements $f_j$ and $h_i$ are mapped to zero
in $\Lambda$. Moreover, it is clear that $g_i$ is mapped to an element
$t$ of $\pi\Lambda$ such that $\pi^3t = 0$ (as $f_i = \pi^3 g_i$ modulo
the ideal generated by the $h$'s). Hence our assumption that
$\text{Ann}_\Lambda(\pi) = \text{Ann}_\Lambda(\pi^2)$ implies that $t = 0$.
Thus we are done if we can prove the statement about smoothness.

\medskip\noindent
Note that $B_\pi \cong A_\pi[v_1, \ldots, v_c]$ because the equations
$g_i = 0$ are implied by $f_i = 0$. Hence $B_\pi$ is smooth over $R$
as $A_\pi$ is smooth over $R$ by the assumption that $\pi$ is strictly
standard in $A$ over $R$, see
Lemma \ref{lemma-elkik}.

\medskip\noindent
Set $B' = R[v_1, \ldots, v_c, w_1, \ldots, w_n]/(g_1, \ldots, g_c)$.
As $g_i = v_i + \sum r_{ji}w_i \bmod \pi$ we see that
$B'/\pi B' = R/\pi R[w_1, \ldots, w_n]$. Hence
$R \to B'$ is smooth of relative dimension $n$ at every
point of $V(\pi)$ by
Algebra, Lemmas
\ref{algebra-lemma-localize-relative-complete-intersection} and
\ref{algebra-lemma-flat-fibre-smooth}
(the first lemma shows it is syntomic at those primes, in particular
flat, whereupon the second lemma shows it is smooth).

\medskip\noindent
Let $\mathfrak q \subset B$ be a prime with $\pi \in \mathfrak q$ and
for some $r \in \mathfrak a$, $r \not \in \mathfrak q$.
Denote $\mathfrak q' = B' \cap \mathfrak q$.
We claim the surjection $B' \to B$ induces an isomorphism of local
rings $(B')_{\mathfrak q'} \to B_\mathfrak q$. This will
conclude the proof of the lemma. Note that $B_\mathfrak q$ is the
quotient of $(B')_{\mathfrak q'}$ by the ideal generated by
$f_{c + j}$, $j = 1, \ldots, m - c$. We observe two things:
first the image of $f_{c + j}$ in $(B')_{\mathfrak q'}$ is
divisible by $\pi^2$ and
second the image of $\pi f_{c + j}$ in $(B')_{\mathfrak q'}$
can be written as $\sum b_{j_1 j_2} f_{c + j_1}f_{c + j_2}$ by
(\ref{equation-star}). Thus we see that the image of each $\pi f_{c + j}$
is contained in the ideal generated by the elements $\pi^2 f_{c + j'}$.
Hence $\pi f_{c + j} = 0$ in $(B')_{\mathfrak q'}$ as this is a
Noetherian local ring, see
Algebra, Lemma \ref{algebra-lemma-intersect-powers-ideal-module-zero}.
As $R \to (B')_{\mathfrak q'}$ is flat we see that
$$
\left(\text{Ann}_R(\pi^2)/\text{Ann}_R(\pi)\right)
\otimes_R (B')_{\mathfrak q'}
=
\text{Ann}_{(B')_{\mathfrak q'}}(\pi^2)/\text{Ann}_{(B')_{\mathfrak q'}}(\pi)
$$
Because $r \in \mathfrak a$ is invertible in
$(B')_{\mathfrak q'}$ we see that this module is zero.
Hence we see that the image of $f_{c + j}$ is zero in
$(B')_{\mathfrak q'}$ as desired.
\end{proof}
